\subsection{双线性型或半双线性型的逆型。}

设 A 是一个环，E 是一个左 A-模，F 是一个右 A-模，\(\Phi\) 是 \(\mathbf{E} \times \mathbf{F}\) 上的一个双线性型。这里假设 \(\Phi\) 的相伴映射（将记作 \(s\) 与 \(d\) ）是双射。于是，积映射 \((s, d)\) 是 \(\mathbf{E} \times \mathbf{F}\) 到 \(\mathbf{F}^* \times \mathbf{E}^*\) 上的一个双射，且经由结构转移，定义了双线性型 \(\hat{\Phi}\) （在 \(\mathbf{F}^* \times \mathbf{E}^*\) 上）。因此它满足

\[
\begin{array}{r l} \hat {\Phi} (y ^ {\prime}, x ^ {\prime}) = & \Phi (s ^ {- 1} (y ^ {\prime}), d ^ {- 1} (x _ {i} ^ {\prime})) \\ = & \left\langle s ^ {- 1} (y ^ {\prime}), x ^ {\prime} \right\rangle = \left\langle d ^ {- 1} (x ^ {\prime}), y ^ {\prime} \right\rangle \qquad (x ^ {\prime} \in \mathrm{E} ^ {*}, y ^ {\prime} \in \mathrm{F} ^ {*}). \end{array}\tag{27}
\]

定义 8. --- 设 \(\Phi\) 是 E × F 上的一个双线性型，它的相伴映射 s 与 d 是双射。由 (27) 定义的 F* × E* 上的双线性型 \(\hat{\Phi}\) 称为 \(\Phi\) 的逆型。

现在设 \(\hat{s}\) 与 \(\hat{d}\) 为左（相应地右）相伴于 \(\hat{\Phi}\) 的、F* 到 E** 与 E* 到 F** 中的线性映射。由于对 \(x' \in \mathrm{E}^*\) 与 \(y' \in \mathrm{F}^*\) ，按定义有

\[
\hat {\Phi} (y ^ {\prime}, x ^ {\prime}) = \left\langle y ^ {\prime}, \hat {d} (x ^ {\prime}) \right\rangle = \left\langle x ^ {\prime}, \hat {s} (y ^ {\prime}) \right\rangle
\]

与 (27) 相比较，便看出 F* 上的线性型 \(\hat{d}(x')\) 等于由 F 的元素 \(d^{-1}(x')\) 所定义的线性型。由此得出，复合映射 \(\hat{d} \circ d\) 是 F 到其双对偶 F** 中的典范映射，且这个映射是双射，因为 \(d\) 与 \(\hat{d}\) （后者经由结构转移）都是双射；于是，若把 F 与 F** 典范地等同起来，就有 \(\hat{d} = d^{-1}\) 。类似地，E 与 E** 典范地等同，而 E 到 E** 的典范映射是 \(\hat{s} \circ s\) ，且有 \(\hat{s} = s^{-1}\) 。由此推出 \(\hat{\Phi}\) 的逆型是 \(\Phi\) 。

现在考虑一个配备了反自同构 J 的环 A、两个左 A-模 E 与 F，以及一个（关于 J 的）右半双线性型 \(\Phi\) （在 E \(\times\) F 上），使得 \(\Phi\) 的相伴映射（将记作 \(s\) 与 \(d\) ）是双射。用 (27) 的第一个等式定义一个映射 \(\hat{\Phi}\) （从 \(\mathbf{F}^{*} \times \mathbf{E}^{*}\) 到 A 中）。由 (24)（第 6 小节），这个映射满足关系式

\[
\hat {\Phi} (y ^ {\prime}, x ^ {\prime}) = \left\langle s ^ {- 1} (y ^ {\prime}), x ^ {\prime} \right\rangle = \left\langle d ^ {- 1} (x ^ {\prime}), y ^ {\prime} \right\rangle^ {J} \quad (x ^ {\prime} \in E ^ {*}, y ^ {\prime} \in F ^ {*}).\tag{28}
\]

映射 \(\hat{\Phi}\) 显然是 \(\mathbf{Z}\) -双线性的；此外，对 A 中的 \(a\) 、\(b\) 、\(x' \in \mathrm{E}^*\) 与 \(y' \in \mathrm{F}^*\) ，鉴于 \(s\) 与 \(d\) 的定义，我们有

\[
\widehat {\Phi} (y ^ {\prime} a, x ^ {\prime} b) = \Phi (a ^ {j} s ^ {- 1} (y ^ {\prime}), b ^ {j ^ {\prime}} d ^ {- 1} (x ^ {\prime})) = a ^ {j} \widehat {\Phi} (y ^ {\prime}, x ^ {\prime}) b;
\]

因此 \(\hat{\Phi}\) 是 \(\mathbf{F}^{*} \times \mathbf{E}^{*}\) 上的一个（关于 J 的）左半双线性型（第 2 小节）。

定义 9. --- 设 \(\Phi\) 是 E × F 上的一个（关于 J 的）右半双线性型，其相伴映射 s 与 d 是双射。F* × E* 上的（关于 J 的）左半双线性型 \(\hat{\Phi}\) 称为 \(\Phi\) 的逆型。

我们把左半双线性型的逆型的定义与研究留给读者。这个逆型是一个右半双线性型。

设 \(\hat{s}\) 与 \(\hat{d}\) 为 \(\hat{\Phi}\) 的相伴映射；由 (26)（第 6 小节），我们有

\[
\hat {\Phi} (y ^ {\prime}, x ^ {\prime}) = \left\langle y ^ {\prime}, \hat {d} (x ^ {\prime}) \right\rangle^ {\mathbf {J}} = \left\langle x ^ {\prime}, \hat {s} (y ^ {\prime}) \right\rangle .\tag{29}
\]

由于 \(s\) 是双射，并由于由 (28) 与 (29) 得出的等式 \(\langle s^{-1}(y'), x'\rangle = \langle y', \hat{d}(x') \rangle^{\mathfrak{J}}\) ，我们导出 \(\hat{d}\) 是双射；于是 \(\hat{d} \circ d\) 是双射。但由 (28) 与 (29) 得出的等式 \(\langle d^{-1}(x'), y' \rangle = \langle y', \hat{d}(x') \rangle\) 表明 \(\hat{d} \circ d\) 是 F 到其双对偶 F** 中的典范映射。同样看出 \(\hat{s}\) 是双射且 \(\hat{s} \circ s\) 是 E 到 E** 中的典范映射。因此，若借助这些典范映射把 E** 与 E、F** 与 F 等同起来，我们就有 \(\hat{s} = s^{-1}\) 、\(\hat{d} = d^{-1}\) ，且 \(\Phi\) 是 \(\hat{\Phi}\) 的逆型。

在同一记号与假设下，设 a 是 A 的中心的一个可逆元素。于是，由 (23)（相应地 (24)），与型 \(a\Phi\) 相伴的映射等于 \(a.d\) 与 \(a.s\) （相应地 \(a^{j'}\) .s），从而是双射。于是由 (27) 推出，\(a\Phi\) 的逆型是 \(a^{-1}\widehat{\Phi}\) （相应地 \((a^{j'})^{-1}\widehat{\Phi}\) ）。

